\subsection{Propiedades de los dominios de definición de las funciones racionales}
\label{IV.20.4}

\begin{env}[20.4.1]
\label{IV.20.4.1}
Sean $X$, $Y$ dos $S$-preesquemas, y sea $\omega$ un pseudo-$S$-morfismo de $X$ en $Y$.
Sea $u$ un $S$-morfismo $U\to Y$ perteneciente a $\omega$, donde $U$ es esquemáticamente denso en $X$, y consideremos el \emph{grafo} $\Gamma_u$, subpreesquema de $U\times_S Y$ \sref[I]{I.5.3.11}, y por tanto subpreesquema de $X\times_S Y$.
Supongamos que este subpreesquema admite una \emph{clausura} \sref[I]{I.9.5.11} en $X\times_S Y$; si $j:\Gamma_u\to X\times_S Y$ es la inyección canónica, así ocurre cuando el $\sh{O}_{X\times_S Y}$-módulo $j_*(\sh{O}_{\Gamma_u})$ es cuasicoherente; de la definición de la clase de equivalencia $\omega$ \sref{IV.20.2.1} se deduce que $j_*(\sh{O}_{\Gamma_u})$ no depende del representante $u$ considerado, por lo que la clausura preesquemática de $\Gamma_u$ queda entonces bien determinada por $\omega$; se dice que este subpreesquema cerrado de $X\times_S Y$ es el \emph{grafo} del pseudo-$S$-morfismo $\omega$, y se denota por $\Gamma_\omega$.
Obsérvese que $\Gamma_\omega$ está definido si existe un morfismo $u:U\to Y$ de la clase $\omega$ tal que $U$ sea \emph{retrocompacto} en $X$ y si, además, $Y$ es cuasiseparado sobre $S$, pues entonces la inyección $j$ es un morfismo cuasicompacto y cuasiseparado (\sref[I]{I.2.2} y \sref[I]{I.7.4}); la primera hipótesis se cumple siempre cuando $X$ es \emph{localmente noetheriano}.
Obsérvese también que, cuando $X$ es \emph{reducido}, lo mismo ocurre con $\Gamma_u$, que es isomorfo a $U$ \sref[I]{I.5.3.11}; entonces $\Gamma_\omega$ existe y no es sino el subpreesquema \emph{reducido} de $X\times_S Y$ cuyo espacio subyacente es la clausura, en $X\times_S Y$, del espacio subyacente a $\Gamma_u$ (\sref[I]{I.5.2.1} y \sref[I]{I.5.2.2}).
Obsérvese finalmente que, si $Y$ es \emph{separado} sobre $S$, $\Gamma_u$ es \emph{cerrado} en $U\times_S Y$ \sref[I]{I.5.4.3}, y por tanto está \emph{inducido por} $\Gamma_\omega$ (cuando este último existe) sobre el abierto $\Gamma_\omega\cap(U\times_S Y)$ \sref[I]{I.9.5.10}; por el contrario, este preesquema inducido es en general distinto de $\Gamma_u$ cuando $Y$ no es separado.
En particular, si $v:X\to Y$ es un $S$-morfismo, el grafo $\Gamma_\omega$ de la clase $\omega$ de $v$ puede ser distinto del grafo $\Gamma_v$.
Por ello, en lo que sigue, cuando se trate del grafo de un pseudo-$S$-morfismo, nos limitaremos al caso en que \emph{$Y$ es separado sobre $S$}.
\end{env}

\begin{env}[20.4.2]
\label{IV.20.4.2}
Supongamos, pues, que $\Gamma_\omega$ está definido y que $Y$ es separado sobre $S$; denotemos por $p$ y $q$ las restricciones a $\Gamma_\omega$ de las proyecciones canónicas
\[
  \xymatrix{
  & \Gamma_\omega \ar[dl]_p \ar[dr]^q & \\
  X & & Y
  }
\]
Entonces, si $U\subset\operatorname{dom}_S(\omega)$, la restricción $p^{-1}(U)\to U$ de $p$ es un isomorfismo \sref[I]{I.5.3.11}; recíprocamente, si $U$ es un abierto de $X$ que tiene esta propiedad, y si $u$ es el isomorfismo inverso de la restricción $p^{-1}(U)\to U$ de $p$, $q\circ u$ es un $S$-morfismo de $U$ en $Y$ que coincide con un morfismo de la clase $\omega$ en $U\cap\operatorname{dom}_S(\omega)$.
Se concluye que $\operatorname{dom}_S(\omega)$ es el \emph{mayor abierto} $U$ de $X$ tal que la restricción $p^{-1}(U)\to U$ de $p$ sea un
\oldpage[IV-4]{245}
isomorfismo.
Sea $v:\operatorname{dom}_S(\omega)\to\Gamma_\omega$ el $S$-morfismo inverso de la restricción $p^{-1}(\operatorname{dom}_S(\omega))\to\operatorname{dom}_S(\omega)$ de $p$; se denota a veces por $p^{-1}$ el pseudo-$S$-morfismo de $X$ en $\Gamma_\omega$, clase de $v$ (cuya aplicación racional asociada \sref{IV.20.2.13}[(ii)] es entonces \emph{birracional}); como $p^{-1}(\operatorname{dom}_S(\omega))$ es el grafo del $S$-morfismo $u=q\circ v:\operatorname{dom}_S(\omega)\to Y$, es esquemáticamente denso en $\Gamma_\omega$ \sref{IV.11.10.3}[(iv)], por lo que $\omega$ puede considerarse como el \emph{compuesto} $q\circ p^{-1}$ \sref{IV.20.3.2}.
Para toda parte $M$ del espacio subyacente a $X$, se pone a veces $\omega(M)=q(p^{-1}(M))$, lo que equivale a considerar $\omega$ como una aplicación de $X$ en $\mathfrak{P}(Y)$ (o, como dicen algunos autores, una «función multiforme»).
Obsérvese que, cuando $x\in\operatorname{dom}_S(\omega)$, $\omega(\{x\})$ es el conjunto $\{u(x)\}$; en general, para un $x\in X$ cualquiera, si se denota por $s$ la imagen de $x$ en $S$ y por $Y_s$ la fibra en el punto $s$ del morfismo estructural $Y\to S$, el conjunto $\omega(\{x\})$ es una parte del preesquema $Y_s$.
\end{env}

\begin{env}[20.4.3]
\label{IV.20.4.3}
En todo el resto de este número nos limitaremos al caso en que $X$ sea \emph{reducido}, de modo que hay entonces identidad entre los pseudo-$S$-morfismos y las $S$-\emph{aplicaciones racionales} \sref{IV.20.2.7}.
Además, el \emph{grafo} de toda $S$-aplicación racional de $X$ en $Y$ está entonces definido \sref{IV.20.4.1}.
\end{env}

\begin{proposition}[20.4.4]
\label{IV.20.4.4}
Sean $X$ un $S$-preesquema localmente íntegro, $Y$ un $S$-preesquema localmente de tipo finito y separado sobre $S$, $\omega$ una $S$-aplicación racional de $X$ en $Y$, y $p:\Gamma_\omega\to X$ la proyección canónica.
Entonces, si $x\in X$ es un punto normal tal que el conjunto $p^{-1}(x)$ contiene un punto aislado, $\omega$ está definida en el punto $x$.
\end{proposition}

\begin{proof}
En efecto, la primera proyección $p_1:X\times_S Y\to X$ es entonces un morfismo separado y localmente de tipo finito, por lo que lo mismo sucede con su restricción $p:\Gamma_\omega\to X$, que es además \emph{birracional}; y, como $\Gamma_\omega$ es reducido y $X$ íntegro, $\Gamma_\omega$ es \emph{íntegro}. Se deduce, pues, de (\textbf{Err}\textsubscript{\textbf{IV}}, 30) que la hipótesis sobre $x$ implica la existencia de un entorno abierto $U$ de $x$ tal que la restricción $p^{-1}(U)\to U$ de $p$ sea un isomorfismo; de ahí la conclusión por \sref{IV.20.4.2}.

Obsérvese que el enunciado \sref{IV.20.4.4} es la formulación original dada por Zariski de su «teorema principal» (salvo que originalmente estaba restringido a los esquemas algebraicos sobre un cuerpo base).
\end{proof}

\begin{proposition}[20.4.5]
\label{IV.20.4.5}
Bajo las hipótesis y notaciones de \sref{IV.20.4.4}, supongamos además que $X$ es normal, y sean $X'$ un preesquema reducido y $f:X'\to X$ un morfismo localmente de tipo finito y universalmente abierto.
Entonces $\omega\circ f$ está definida, y se tiene $\operatorname{dom}_S(\omega\circ f)=f^{-1}(\operatorname{dom}_S(\omega))$ (en otras palabras, si $x'\in X'$ y $x=f(x')$, para que $\omega$ esté definida en el punto $x$ es necesario y suficiente que $\omega\circ f$ lo esté en el punto $x'$).
\end{proposition}

\begin{proof}
Como $X'$ es reducido y $f^{-1}(\operatorname{dom}_S(\omega))$ es denso en todo $X'$ en virtud de \sref{IV.2.4.11}, el compuesto $\omega\circ f$ está definido; para demostrar que, cuando $\omega\circ f$ está definida en el punto $x'$, $\omega$ lo está en el punto $x$, podemos sustituir evidentemente $X'$ por un entorno abierto de $x'$ y suponer, por tanto, que $\omega\circ f$ está \emph{definida en todas partes}; además, como $f(X')$ es abierto en $X$, podemos suponer que $f$ es \emph{sobreyectivo}.
En virtud de la hipótesis sobre $f$, el morfismo $f_{(Y)}:X'\times_S Y\to X\times_S Y$ es entonces abierto, por lo que $\overline{f_{(Y)}^{-1}(M)}=f_{(Y)}^{-1}(\overline{M})$ para toda parte $M$ de $X\times_S Y$; tomando como $M$ el conjunto $\Gamma_u$, donde $u:\operatorname{dom}_S(\omega)\to Y$ es la restricción de $\omega$ a $\operatorname{dom}_S(\omega)$, se deduce de la relación anterior y de \sref[I]{I.3.3.12} que el conjunto subyacente a $\Gamma_{\omega\circ f}$ es igual a $f_{(Y)}^{-1}(\Gamma_\omega)$; como se sabe
\oldpage[IV-4]{246}
que $\Gamma_{\omega\circ f}$ es un preesquema reducido \sref{IV.20.4.1}, se ve que el $X'$-preesquema $\Gamma_{\omega\circ f}$ es igual a $(\Gamma_\omega\times_X X')_\mathrm{red}$.
Pero, como por hipótesis el morfismo compuesto $\Gamma_{\omega\circ f}\to\Gamma_\omega\times_X X'\xrightarrow{p^{(X')}}X'$ es un isomorfismo, $p^{(X')}$ es necesariamente \emph{radicial}; como $f$ es sobreyectivo, lo mismo sucede con $p:\Gamma_\omega\to X$ \sref[I]{I.3.4.8}, por lo que $p^{-1}(x)$ es un conjunto reducido a un punto \sref[I]{I.3.6.4}; se deduce entonces de \sref{IV.20.4.4} que $\omega$ está definida en el punto $x$.
\end{proof}

La proposición siguiente da un \emph{criterio valuativo} para que una aplicación racional esté definida en un punto:

\begin{proposition}[20.4.6]
\label{IV.20.4.6}
Sean $S$ un preesquema, y $X$, $Y$ dos $S$-preesquemas; se supone que $X$ es localmente noetheriano y que $Y$ es localmente de tipo finito sobre $S$.
Sean $U$ un abierto denso de $X$, $h:U\to Y$ un $S$-morfismo, $x\in X-U$ un punto normal de $X$, y $h'_x:\Spec(k(x))\to Y$ un $S$-morfismo.
Para que $h$ pueda prolongarse a un $S$-morfismo $h':U'\to Y$, donde $U'$ es un abierto de $X$ que contiene $U$ y $x$, y donde el morfismo compuesto $\Spec(k(x))\to U'\xrightarrow{h'}Y$ es el $S$-morfismo dado $h'_x$, es necesario y suficiente que se cumpla la condición siguiente:
\begin{quote}
\emph{(P) Para todo espectro $S_1$ de un anillo de valoración discreta, con punto cerrado $a$ y punto genérico $b$, y todo morfismo $u:S_1\to X$ tal que $u(a)=x$ y $u(b)\in U$, el morfismo compuesto $h_1=h\circ(u|\{b\}):\{b\}=u^{-1}(U)\to Y$ se prolonga a un morfismo $h'_1:S_1\to Y$ tal que el diagrama}
\[
  \xymatrix{
  \Spec(k(a)) \ar[r] & S_1 \ar[d]^{h'_1} \\
  \Spec(k(x)) \ar[r]_{h'_x} & Y
  }
\]
\emph{sea conmutativo.}
\end{quote}
Además, si se cumple esta condición y si $h'':U''\to Y$ es un morfismo que satisface las mismas condiciones que $h'$, existe un conjunto abierto que contiene $U\cup\{x\}$ y en el cual $h'$ y $h''$ coinciden.
\end{proposition}

\begin{proof}
Demostremos primero la última afirmación; podemos suponer que $h'$ y $h''$ están definidos en el mismo abierto $U_0\supset U\cup\{x\}$.
El subpreesquema $Z$ de coincidencia de $h'$ y $h''$ \sref{IV.17.4.5} contiene $U$ y $x$, por lo que existe un entorno abierto $V$ de $x$ en $U_0$ tal que el subpreesquema inducido por $Z$ sobre el abierto $Z\cap V$ sea un subpreesquema \emph{cerrado} del preesquema inducido por $X$ sobre $V$; como este preesquema $Z\cap V$ contiene el subpreesquema inducido por $V$ sobre el abierto $U\cap V$, y este último es esquemáticamente denso en $V$, $Z\cap V$ es necesariamente igual a $V$ \sref{IV.20.3.8.8}, lo que demuestra que $h'$ y $h''$ coinciden en $U\cup V$.

Siendo evidente la necesidad de la condición del enunciado, demostremos que es suficiente.
En virtud de la correspondencia biyectiva entre los $S$-morfismos de $X$ en $Y$ y las $X$-secciones de $X\times_S Y$ \sref[I]{I.3.3.15}, podemos limitarnos al caso en que $S=X$ y en que $h$ es, por tanto, una $U$-sección de $Y$; obsérvese que entonces $Y$ es localmente noetheriano, y podemos evidentemente (puesto que $X$ es localmente íntegro y localmente noetheriano) suponer que $X=S$ es irreducible.
Además, teniendo en cuenta \sref{IV.20.3.7}, podemos suponer que $X=\Spec(A)$, donde $A$ es un anillo \emph{local} noetheriano, íntegro e integralmente cerrado.

\oldpage[IV-4]{247}
Pongamos $y=h'_x(x)$. Para todo $x'\in U$, $y$ es una especialización de $h(x')$.\footnote{La transcripción francesa invierte aquí la relación de especialización; el testigo inglés mantenido y las igualdades $h(x')=h'_1(b)$, $y=h'_1(a)$ muestran que $y$ es la especialización de $h(x')$.}
En efecto, existe un espectro $S_1$ de un anillo de valoración discreta y un morfismo $u:S_1\to X$ tal que $u(a)=x$ y $u(b)=x'$ \sref[II]{II.7.1.9}.
Aplicando la condición del enunciado, se obtiene inmediatamente nuestra afirmación, pues $h(x')=h_1(b)=h'_1(b)$ e $y=h'_1(a)$.
Si $Y'$ es un entorno abierto afín de $y$, se tiene, por tanto, $h(X\cap U)\subset Y'$; podemos sustituir entonces $Y$ por $Y'$, es decir, suponer que $Y$ es \emph{afín}, y por tanto separado sobre $X$.
Sea $\omega$ la $X$-sección racional de $Y$ a la que pertenece $h$, de modo que su grafo tiene aquí por conjunto subyacente la clausura de $h(U)$ en $Y$.
Como $Y$ es separado sobre $X$, podemos aplicar \sref{IV.20.4.4} a $\omega$: bastará demostrar que, si $p:\Gamma_\omega\to X$ es la proyección canónica, $p^{-1}(x)$ se reduce al único punto $y$ y que $h'_x(x)=y$.
En efecto, por \sref{IV.20.4.4}, $h$ se prolongará a una sección $h'$ sobre un abierto $U'$ que contiene $U\cup\{x\}$, tal que $h'(x)=y$; y, como no existe más que un $X$-morfismo $\Spec(k(x))\to Y$ que transforma $x$ en $y$, esto demostrará la identidad de $h'_x$ con el compuesto de $h'$ y $\Spec(k(x))\to U'$.

Como, para todo $x'\in X\cap U$, $y$ es una especialización de $h(x')$, se tiene $y\in p^{-1}(x)$.
Supongamos ahora que $z$ es un punto cualquiera de $p^{-1}(x)$.
Como $\Gamma_\omega$ es la clausura de $h(U)$ y $h(U)$ está formado por puntos adherentes a $h(\xi)$, donde $\xi$ es el punto genérico de $X$, $\Gamma_\omega$ es la clausura de $h(\xi)$ en $Y$.
Se toma entonces un espectro $S_1$ de un anillo de valoración discreta y un morfismo $v:S_1\to Y$ tales que $v(a)=z$ y $v(b)=h(\xi)$ \sref[II]{II.7.1.9}, y se pone $u=p\circ v$, de modo que $u(a)=x$ y $u(b)=\xi$.
Aplicando a $u$ la condición del enunciado, se obtiene un morfismo $h'_1:S_1\to Y$ tal que $h'_1(a)=y$ y $h'_1(b)=h(\xi)$; pero esto implica $z=y$ en virtud de \sref[II]{II.7.2.3}, puesto que $Y$ es separado sobre $X$ y, por tanto, $v$ y $h'_1$ deben coincidir, ya que son iguales en el punto $b$.
C.Q.F.D.
\end{proof}

\begin{corollary}[20.4.7]
\label{IV.20.4.7}
Bajo las hipótesis sobre $S$, $X$, $Y$, $U$ y $h$ de \sref{IV.20.4.6}, sea $E$ una parte de $X-U$ tal que $X$ sea normal en todo punto de $E$, y, para cada $x\in E$, sea $h'_x:\Spec(k(x))\to Y$ un $S$-morfismo tal que se cumpla la condición~(P) de \sref{IV.20.4.6}.
Supongamos además que la unión $F$ de los grafos de los $h'_x$ (identificados con partes de $X\times_S Y$), para $x\in E$, está contenida en la unión de un número finito de abiertos $V_i$ de $X\times_S Y$ que son \emph{separados} sobre $X$ (condición que se cumple automáticamente si $Y$ es \emph{separado} sobre $S$, o si $X$ es \emph{cuasicompacto} e $Y$ de \emph{tipo finito} sobre $S$).
Entonces existe un $S$-morfismo $h':U'\to Y$, donde $U'$ es un abierto de $X$ que contiene $U\cup E$, tal que, para todo $x\in E$, el compuesto
\[
  \Spec(k(x))\to U'\xrightarrow{h'}Y
\]
sea igual a $h'_x$.
\end{corollary}

\begin{proof}
Obsérvese primero que, en \sref{IV.20.4.6}, si se supone que $Y$ es \emph{separado}, existe un mayor abierto $U_0\supset U$, igual al dominio de la $S$-aplicación racional correspondiente a $h$, en el que $h$ puede prolongarse, y este prolongamiento es único \sref[I]{I.7.2.2}; de ahí la conclusión en este caso, en virtud de \sref{IV.20.4.6}.
En el caso general, sea $E_i$ el conjunto de los $x\in E$ tales que $(x,h'_x(x))\in V_i$.
En virtud de \sref{IV.20.4.6}, para todo $x\in E$ existe un prolongamiento $h^{(x)}$ de $h$ a $U\cup W^{(x)}$, donde $W^{(x)}$ es un entorno abierto de $x$ en $X$ tal que el grafo de $h^{(x)}|W^{(x)}$ está contenido en \emph{todos} los $V_i$ tales que $x\in E_i$.
Como $V_i$ es separado sobre $X$ y $X$ es reducido, para dos puntos $x'$, $x''$ de $E_i$, $h^{(x')}$ y $h^{(x'')}$ coinciden en $W^{(x')}\cap W^{(x'')}$, puesto que
\oldpage[IV-4]{248}
coinciden en el abierto denso en todas partes $U\cap W^{(x')}\cap W^{(x'')}$ de este conjunto.
Existe, pues, un morfismo $h_i:U\cup U_i\to Y$ que prolonga $h$ a un abierto $U\cup U_i$ que contiene $E_i$; además, para todo par de índices $i$, $j$, los grafos de las restricciones $h_i|(U_i\cap U_j)$ y $h_j|(U_i\cap U_j)$ están contenidos en $V_i\cap V_j$; como $V_i\cap V_j$ es separado sobre $X$ y los morfismos anteriores coinciden en un abierto denso en todas partes $U\cap U_i\cap U_j$ de $U_i\cap U_j$, son iguales.
El morfismo $h'$, igual a $h$ en $U$ y a $h_i$ en cada uno de los $U_i$, responde a la cuestión.
\end{proof}

\begin{remark}[20.4.8]
\label{IV.20.4.8}
Cuando $E=X-U$, es claro que, si para todo abierto afín $T$ de $X$ existe un $S$-morfismo $h'_T:T\to Y$ que prolonga $h|(U\cap T)$ y tal que el compuesto $\Spec(k(x))\to T\xrightarrow{h'_T}Y$ es igual a $h'_x$ para todo $x\in T-(U\cap T)$, entonces los $h'_T$ son las restricciones de un $S$-morfismo $h':X\to Y$ definido en todas partes y que prolonga $h$, en virtud de la afirmación de unicidad de \sref{IV.20.4.6}.
Para demostrar la existencia de $h'$, nos reducimos, pues, al caso en que $X$ es \emph{afín}, y basta entonces que el conjunto $F$, unión de los grafos de los $h'_x$, sea \emph{cuasicompacto} en $X\times_S Y$ para poder aplicar \sref{IV.20.4.7}.
Así sucede cuando los $h'_x$ son de la forma $\Spec(k(x))\to Z\xrightarrow{h_0}Y$, donde $Z$ es un subpreesquema cerrado de $X$ cuyo espacio subyacente es $X-U$, y $h_0$ es un $S$-morfismo.
\end{remark}

\begin{corollary}[20.4.9]
\label{IV.20.4.9}
Sean $S$ un preesquema localmente noetheriano, $X$ un preesquema localmente noetheriano, $f:X\to S$ un morfismo plano y $g:Y\to S$ un morfismo localmente de tipo finito.
Sean $U$ un abierto denso de $S$, $h:f^{-1}(U)\to Y$ un $S$-morfismo, $T=S-U$, $Z$ el subpreesquema cerrado reducido de $X$ cuyo espacio subyacente es $f^{-1}(T)$, y $h_0:Z\to Y$ un $S$-morfismo.
Supongamos que $X$ es normal en todos los puntos de $Z$.
Para que exista un $S$-morfismo (necesariamente único) $h':X\to Y$ que prolongue $h$ y $h_0$, es necesario y suficiente que se cumpla la condición siguiente:

\emph{Para todo espectro $S_1$ de un anillo de valoración discreta, con punto cerrado $a$ y punto genérico $b$, y todo morfismo $u:S_1\to S$ tal que $u(a)\in T$ y $u(b)\in U$, existe un $S_1$-morfismo $h'_1:X_{(S_1)}\to Y_{(S_1)}$ que prolonga $h_{(S_1)}:f_{(S_1)}^{-1}(b)\to Y_{(S_1)}$ y es tal que el diagrama}
\[
  \xymatrix{
  \Spec(k(z_1)) \ar[r] \ar[d] & Z_{(S_1)} \ar[d]^{(h_0)_{(S_1)}} \\
  X_{(S_1)} \ar[r]_{h'_1} & Y_{(S_1)}
  }
\]
\emph{sea conmutativo para todo $z_1\in Z_{(S_1)}$.}

La hipótesis de que $f$ es plano implica, en efecto, que $f^{-1}(U)$ es denso en $X$ \sref{IV.2.3.10}, y basta entonces aplicar \sref{IV.20.4.8}.
\end{corollary}

\begin{corollary}[20.4.10]
\label{IV.20.4.10}
Bajo las hipótesis de \sref{IV.20.4.6}, supongamos además que $Y$ es separado y localmente cuasifinito sobre $S$.
Sean $U$ un abierto denso de $X$, $h:U\to Y$ un $S$-morfismo, $\omega$ la $S$-aplicación racional correspondiente, y $x\in X-U$ un punto normal de $X$.
Para que $\omega$ esté definida en el punto $x$, es necesario y suficiente que se cumpla la condición siguiente: existe un espectro $S_1$ de un anillo de valoración discreta, con punto cerrado $a$ y punto genérico $b$, y un morfismo $u:S_1\to X$
\oldpage[IV-4]{249}
tales que $u(a)=x$, $u(b)\in U$ y la $S$-aplicación racional $\omega\circ u$ está definida en todas partes.

En efecto, por hipótesis todas las fibras del morfismo de proyección $X\times_S Y\to X$ están formadas por puntos aislados, y para aplicar \sref{IV.20.4.4} basta saber que la fibra $p^{-1}(x)$ no es vacía en $\Gamma_\omega$.
Ahora bien, si $h_1$ es el único morfismo de la clase $\omega\circ u$, el único punto de $X\times_S Y$ situado sobre $x$ y sobre $h_1(a)$ pertenece a $\Gamma_\omega$, de donde se obtiene la conclusión.
\end{corollary}

\begin{proposition}[20.4.11]
\label{IV.20.4.11}
Sean $X$ un $S$-preesquema localmente noetheriano,\footnote{La transcripción francesa omite aquí la estructura de $S$-preesquema sobre $X$, necesaria para que los morfismos del enunciado sean $S$-morfismos y usada explícitamente en la demostración.} $Y$ un $S$-preesquema afín sobre $S$, $U$ un abierto de $X$, y $Z=X-U$.
Supongamos que $\operatorname{prof}_Z(\sh{O}_X)\geq 2$ \sref{IV.5.10.1}; entonces todo $S$-morfismo $f:U\to Y$ se prolonga de manera única a un $S$-morfismo de $X$ en $Y$.
\end{proposition}

\begin{proof}
Podemos limitarnos al caso en que $S$ y $X$ sean afines y (en virtud de \sref[I]{I.3.3.14}) al caso en que $S=X$; se tiene, pues, $X=\Spec(A)$, $Y=\Spec(B)$, siendo $B$ una $A$-álgebra.\footnote{La transcripción francesa y el testigo inglés añaden aquí «de tipo finito», hipótesis que no figura en el enunciado ni se necesita: toda $A$-álgebra es cociente de un álgebra de polinomios en una familia posiblemente infinita de indeterminadas.}
Como $B$ es cociente de un álgebra de polinomios $A[(T_\lambda)]_{\lambda\in L}$, $Y$ es un subpreesquema cerrado de $Y'=X[(T_\lambda)]_{\lambda\in L}$.
Por otra parte, la hipótesis sobre $Z$ implica que $U$ es esquemáticamente denso en $X$ en virtud de \sref{IV.20.2.13}[(iv)] y \sref{IV.5.10.2}.
Si se demuestra que todo $X$-morfismo $u:U\to Y$ se prolonga de manera única a un $X$-morfismo $v':X\to Y'$, se deducirá que $v'$ se factoriza como $X\xrightarrow{v}Y\to Y'$; en efecto, el subpreesquema $v'^{-1}(Y)$ es cerrado y contiene $U$ \sref[I]{I.4.4.1}, por lo que es idéntico a $X$ \sref{IV.20.3.8.8}.
En estas condiciones, $v$ será el único $S$-morfismo de $X$ en $Y$ que prolonga $u$.
Podemos, pues, limitarnos al caso en que $Y=Y'$.
Pero entonces hay una correspondencia biyectiva entre los $X$-morfismos de un abierto $U\subset X$ en $Y'$ y las familias $(s_\lambda)_{\lambda\in L}$ de secciones de $\sh{O}_X$ sobre $U$ \sref[II]{II.1.7.9}; la conclusión se deduce, pues, de \sref{IV.5.10.5}.
\end{proof}

\begin{corollary}[20.4.12]
\label{IV.20.4.12}
Sean $X$ un $S$-preesquema localmente noetheriano, reducido y que satisface la condición $(S_2)$ \sref{IV.5.7.2} (por ejemplo, un $S$-preesquema localmente noetheriano y normal \sref{IV.5.8.6}), $Y$ un $S$-preesquema afín sobre $S$, y $f$ una $S$-aplicación racional de $X$ en $Y$; entonces toda componente irreducible de $X-\operatorname{dom}(f)$ es de codimensión $1$ en $X$.
\end{corollary}

\begin{proof}
Esto equivale a decir que, si $Z^{(2)}$ es el conjunto de los $x\in X$ tales que $\dim(\sh{O}_{X,x})\geq 2$, entonces, para toda parte cerrada $Z\subset Z^{(2)}$ de $X$, todo $S$-morfismo de $X-Z$ en $Y$ se prolonga a un $S$-morfismo de $X$ en $Y$; ahora bien, esto se deduce de la hipótesis sobre $X$ \sref{IV.5.7.2} y de \sref{IV.20.4.11}.
\end{proof}
